\honest{The Importance of Saying ``Oops''}{Eliezer Yudkowsky}{2007}

I have just read a history of Enron's collapse, \textsc{The Smartest Guys in the Room}, and I give it my award for ``Least Appropriate Book Title.''

Enron's executives never admitted a large mistake. When a disaster forced a change of policy, they said, ``Too bad that didn't work out,'' and asked how to hide the problem on the balance sheet. I put that line in their mouths. They never said ``I've been stupid''; there was no watershed moment. After the bankruptcy, Jeff Skilling declined his lawyers' advice to take the Fifth Amendment and told Congress that Enron had been a great company. \nb{The invented line shows people who knew about the problem and hid it. Skilling had been convicted of fraud more than a year before the post, which does not mention it.}

My thesis: every improvement is a change, and if you admit only small mistakes you will make only small changes. A big change needs a big admission. My evidence is Enron and my own case. \nb{In the post's own case, small concessions did add up to a full admission; the complaint is that they were slow.}

I grew up on Heinlein and Feynman, and learned what I call Traditional Rationality: make bold theories that can be proved false, give up your ideas when the evidence goes against them, argue fairly, do not deceive yourself. I call these ``fuzzy verbalisms.'' Traditional Rationality teaches you to give in to the evidence eventually, which is no small thing, since it separates science from religion. But it does not teach speed: using evidence efficiently, so that the least contrary evidence needed is enough to overturn a cherished belief.

Then I made a large mistake, and switched to Bayes. What was the mistake? ``well, it's a long story.'' I do not tell it. This is an essay about admitting large mistakes fully.

Looking back, I had admitted my mistake in small, grudging pieces. I could have moved much faster by simply shouting ``OOPS!''

So admit a large mistake all at once. It is painful, and it can change your whole life. Do not be proud of admitting errors; it is better to be right the first time. ``Even hedonically, it is better to take one large loss than many small ones.'' ``The alternative is Enron.''

Since then I have watched others concede each millimeter of ground, confessing a local mistake where a global one is due. What they could fix at once, voluntarily, they turn into small patches they must be argued into. After one mistake they never say ``I've been a fool,'' but ``I was right in principle,'' ``It could have worked,'' or that they still want the true essence of whatever they are attached to. Defending their pride now, they ensure they will make the same mistake again.

``Better to swallow the entire bitter pill in one terrible gulp.''
